\section{The short-basis and common-curve constructions}
\label{p14:comp:rigidity}

The notation is that of Section~\ref{p13:rig:section}.  This section
supplies the complete local estimates, basis extraction, and algebraic
construction used by the main rank argument.  The final
degree-uniform point count and the nonaligned contradiction are proved
in that section of the article.

\subsection{Local square relations and defective positions}

For primes larger than the window length, a local square relation
can use only positions having no odd valuation at any such prime.
The next elementary argument bounds those positions without an
admissibility assumption.

\begin{lemma}[A local square threshold]
\label{p13:rig:local-square}
Let $B\ge2$, let $I\subset[-B,B]\cap\Z$ have positive even
cardinality $m$, and let $x$ be a positive integer.  If
$x>(B+1)^{12m}$, then $\prod_{a\in I}(x+a)$ is not an integer
square.
\end{lemma}

\begin{proof}
Write $m=2n$ and
\[
 P(T)=\prod_{a\in I}(T+a),\qquad
 \sqrt{P(T)}=T^n\Psi(T^{-1}),\qquad
 \Psi(z)=\prod_{a\in I}(1+az)^{1/2}
                  =\sum_{j\ge0}q_jz^j.
\]
The binomial branches have constant term one.  Their coefficients
give $q_j\in4^{-j}\Z$.  With $R=2(B+1)$, Cauchy's estimate on
$|z|=R^{-1}$ gives $|q_j|\le2^mR^j$.  The polynomial part
$Q(T)=\sum_{j=0}^nq_jT^{n-j}$ has denominator dividing $M=4^n$,
and, for $x\ge2R$,
\[
 |\sqrt{P(x)}-Q(x)|\le\frac{2^{m+1}R^{n+1}}x.
\]
The assumed threshold makes this less than $1/M$, since
\[
 2^{m+1}R^{n+1}M
 =2^{5n+2}(B+1)^{n+1}
 \le(B+1)^{6n+3}\le(B+1)^{9n}.
\]
If $\sqrt{P(x)}$ were an integer, it would therefore equal $Q(x)$.

However, $E=M^2(P-Q^2)$ is a nonzero integer polynomial: the
distinct linear factors make $P$ squarefree.  Its coefficient
height satisfies
\[
 \|P\|_1\le(B+1)^{2n},\qquad
 \|Q\|_1\le2^{3n+1}(B+1)^n\le(B+1)^{5n},
 \qquad H(E)\le(B+1)^{15n}.
\]
Cauchy's root bound for the integer polynomial $E$ would give
$x\le1+H(E)\le(B+1)^{16n}$, contrary to
$x>(B+1)^{24n}$.
\end{proof}

\begin{lemma}[Few defective positions in every critical window]
\label{p13:rig:defective}
The number of $n\in V_x$ for which
\[
             \prod_{\substack{p>10B\\v_p(n)\ {\rm odd}}}p=1
\]
is $O_C(B/\log B)$, uniformly over all critical windows.
\end{lemma}

\begin{proof}
We use the elementary bound $\pi(t)\ll t/\log t$.  For
completeness, the product of the primes in $(u,2u]$ divides
$\binom{2u}{u}\le4^u$ for integer $u$; dyadic summation gives
$\vartheta(t)=\sum_{p\le t}\log p\ll t$.  Splitting at
$\sqrt t$ gives the stated bound for $\pi(t)$.

Let $D$ be the number of defective positions and
$s=\pi(10B)$.  Their square classes lie in an $s$-dimensional
space.  Set
\[
              k=\left\lfloor
                    \frac{\log N}{48\log(B+1)}\right\rfloor.
\]
For sufficiently large $N$, $k\asymp_C B/\log B$ and $k\ge1$.
Lemma~\ref{p13:rig:local-square} forbids a nonempty even local
relation of size at most $2k$, because
$24k\log(B+1)\le\frac12\log N<\log x$.
Two distinct $k$-subsets of the defective positions therefore
have distinct square classes: their symmetric difference would
be such a relation.  If $D\ge k$, this gives
$\binom Dk\le2^s$.  When $D\ge2k$, the inequality
$\binom Dk\ge(D/k)^k$ implies $D\le k2^{s/k}=O_C(k)$;
the case $D<2k$ has the same bound.
\end{proof}

We will also need the elementary prime harmonic estimate
\begin{equation}
 \sum_{B^a<p\le10B}\frac1p
       =\log(1/a)+O_a(1/\log B),\qquad 0<a<1\text{ fixed}.
 \label{p13:rig:mertens}
\end{equation}
Here is a derivation recording its error term.  Chebyshev's bound
above implies $\psi(t)=\sum_{p^j\le t}\log p\ll t$.
The identity
$\log(n!)=\sum_{d\le n}\Lambda(d)\lfloor n/d\rfloor$
then gives $\sum_{d\le t}\Lambda(d)/d=\log t+O(1)$.
The contribution of prime powers with exponent at least two is
bounded, so $A(t)=\sum_{p\le t}(\log p)/p=\log t+O(1)$.
Partial summation yields
\[
 \sum_{p\le t}\frac1p
 =\frac{A(t)}{\log t}+\int_2^t\frac{A(u)}{u(\log u)^2}\,du
 =\log\log t+c_0+O(1/\log t).
\]
The difference at $10B$ and $B^a$ proves
\eqref{p13:rig:mertens}.

\subsection{Large primes turn intersections into components}

Construct a graph on the disjoint formal copies of $V_x,V_y$.
A prime $p>10B$ has at most one occurrence of odd valuation
in either block.  If it has two occurrences, join their vertices;
if it has only one, mark that vertex as forced to zero.
The corresponding binary equation says that the endpoints of
an edge have equal coordinates, or that a marked coordinate is
zero.  A component is free when it contains no marked vertex.

\begin{lemma}[Free components]
\label{p13:rig:components}
Let $c$ be the number of free components containing an edge.
Then
\[
                         c=d_{xy}+O_C(B/\log B).
\]
Each such component meets both blocks and is selected in its
entirety or omitted in any solution to the large-prime equations.
\end{lemma}

\begin{proof}
The isolated unmarked vertices are precisely the defective
positions.  Their total number $D_0$ is $O_C(B/\log B)$ by
Lemma~\ref{p13:rig:defective}.  The kernel of the large-prime
equations has dimension $c+D_0$.  Imposing the remaining
prime equations and the two block parities costs at most
$\pi(10B)+2$ dimensions.  If $\mathcal R$ is the space of
complete relations and $k_x=\dim\ker\phi_x$, then the map
taking a relation to its common class gives an exact sequence
with kernel $\ker\phi_x\oplus\ker\phi_y$.  Hence
\[
 c+D_0-\pi(10B)-2\le\dim\mathcal R=k_x+k_y+d_{xy}
                                      \le c+D_0.
\]
A local relation can use only defective positions, so
$k_x+k_y\le D_0$.  These inequalities prove the assertion.
\end{proof}


\subsection{Proof of Theorem~\ref{p13:rig:short-basis}}
\label{p14:comp:rigidity:short-basis}

\begin{proof}
\emph{A graph estimate.}
A multigraph with $n$ vertices and $e/n\to\infty$ has a cycle
of length $O(1+\log(n+1)/\log(e/n))$.  Loops and parallel
edges give cycles of length one and two.  Otherwise delete
vertices of degree less than $t=\lfloor e/(2n)\rfloor$.
Not all vertices can disappear, since fewer than $nt\le e/2$
edges would then be deleted.  The remaining graph has minimum
degree at least $t$.  If its girth exceeds $2r+1$, a breadth
first search contains at least
$1+t\sum_{j<r}(t-1)^j$ vertices.  This bounds $r$ and proves
the estimate.

Consequently, when $n\ll B/\log B$ and $e\le B$, the binary
cycle space contains at least $e-O(B/\sqrt{\log B})$
independent cycles of length $O(\log B/\log\log B)$.
Indeed, while at least $B/\sqrt{\log B}$ edges remain,
record a short cycle and delete just one of its edges.
The recorded cycles are independent: the first deleted edge
belongs to no later cycle, and the same argument applies
successively to all deleted edges.

\emph{Medium-prime columns.}
Start with the free nontrivial components of
Lemma~\ref{p13:rig:components}, and discard those with more
than $M=\lfloor\log B\rfloor$ vertices.  This loses at most
$2B/(M+1)$ components.  The remaining number is
$d_{xy}+O_C(B/\log B)$.  Put $z=B^a$.  To each retained
component attach its parity column on the primes $z<p\le10B$,
and let $w_j$ be that column's Hamming weight.  Each nonzero
entry consumes a divisible position, whence
\[
 \sum_j w_j
 \le\sum_{z<p\le10B}\#\{n\in V_x\sqcup V_y:p\mid n\}
 \le2B\sum_{z<p\le10B}\frac1p+2\pi(10B)
 \le2B\log(1/a)+O_a(B/\log B).
\]
Discard only columns of weight at least three.  Their number
is at most one third of this sum, so at least
\[
                 e\ge d_{xy}-\frac23B\log(1/a)
                                   -O_{a,C}(B/\log B)
\]
columns remain.  There is no factor $M$ in this count: it
bounds the number of discarded columns, not their physical size.

Use the primes $z<p\le10B$ and one extra vertex $\infty$
as vertices of a multigraph.  A column of weight two joins its
two primes; a column of weight one joins its prime to $\infty$;
a zero column is a loop at $\infty$.  Its binary incidence
matrix has the same kernel as the parity matrix, because the
row at $\infty$ is the sum of the other rows.  The graph has
$n\le\pi(10B)+1\ll B/\log B$ vertices and $e\le B$ edges.
The preceding graph argument therefore supplies a space
$\mathcal U$ of dimension at least
$e-O(B/\sqrt{\log B})$ with a basis of short cycles.
Each basis vector uses
$O(\log^2 B/\log\log B)$ physical positions.

\emph{Small primes and parity.}
All large and medium primes already cancel on $\mathcal U$.
The primes at most $z$ and the two block parities impose at
most $s=\pi(z)+2\ll_a B^a/\log B+2$ further equations on its
cycle coordinates.  Their kernel is $\mathcal K_a$; since
$s=o(B/\sqrt{\log B})$, it has the dimension in
\eqref{p13:rig:short-dimension}.  A fundamental kernel basis
for a matrix with at most $s$ rows uses at most $s+1$ columns
per vector: express each nonpivot column in terms of pivot
columns.  Thus each resulting vector uses at most $s+1$
short cycles, giving \eqref{p13:rig:short-support}.

Finally, the original free nontrivial components are disjoint
and meet both blocks.  Every nonzero selection of them contains
a component carrying a large prime which occurs once in each
block and in no other component.  Its common square class is
therefore nonzero.  This also proves injectivity of both physical
projections and of the common-class map.  The imposed parity
equations complete the proof.
\end{proof}

\subsection{Two relations produce a polynomial of controlled height}

The short basis will have support $B^{a+o(1)}$.  We therefore
need an auxiliary polynomial whose degree and logarithmic
height grow linearly with the support, and a simultaneous
construction whose threshold is quadratic in that support.
The following elementary factor bound will keep the height
under control when common polynomial factors are removed.

\begin{lemma}[Height of a bivariate factor]
\label{p13:rig:factor-height}
Let $P\in\Z[X,Y]\setminus\{0\}$, with partial degrees
$d_X,d_Y$.  Every primitive integer factor $Q$ of $P$, and
the corresponding integer quotient, satisfy
\[
 H(Q)\le2^{d_X+d_Y}\|P\|_2
       \le2^{d_X+d_Y}\sqrt{(d_X+1)(d_Y+1)}\,H(P).
\]
\end{lemma}

\begin{proof}
The multivariate Mahler measure $M(P)$ is multiplicative,
at most $\|P\|_2$, and at least one for a nonzero integer
polynomial.  Jensen's formula in one variable, followed by
the leading coefficient in that variable, proves the last
assertion by iteration.  The one-variable coefficient bound
also iterates: integrating first in $X$ and then in $Y$ gives
\[
 |[X^iY^j]P|\le\binom{d_X}{i}\binom{d_Y}{j}M(P).
\]
Thus $H(Q)\le2^{\deg_XQ+\deg_YQ}M(Q)$; multiplicativity
and $M(P/Q)\ge1$ prove the result.  The quotient has the same
bound, with the integer content handled separately.
\end{proof}

\begin{proposition}[Two complete witnesses]
\label{p13:rig:two-witnesses}
There is an absolute constant $c>0$ with the following property.
Let $u,v$ be independent formal words, even in each block,
with offsets in $[-B,B]$.  For $w=u,v,u+v$, write
$|I_w|=2r_w$, $|J_w|=2s_w$, and assume
$\max(r_w,s_w)\le q$, where $1\le q\le B$.
Suppose that $F_u(x,y)$ and $F_v(x,y)$ are integer squares,
with $x,y\in[N,2N)$.  If $N>(B+1)^{cq}$, there is a
nonzero $R\in\Z[X,Y]$ such that
\begin{equation}
 R(x,y)=0,\qquad
 \deg_XR,\deg_YR\le12q,\qquad \deg R\le24q,
 \qquad H(R)\le(B+1)^{cq}.
 \label{p13:rig:two-polynomial}
\end{equation}
Moreover, $R$ is the norm of a biquadratic auxiliary function
whose four polynomial coefficients have no common nonconstant
factor.
\end{proposition}

\begin{proof}
Set $F_1=F_u$, $F_2=F_v$, $G=\gcd(F_1,F_2)$, with monic
gcd, and $F_3=F_{u+v}=F_1F_2/G^2$.
Formal independence makes
\[
 K=\mathbb Q(X,Y)(W_1,W_2),\qquad W_i^2=F_i,
\]
a degree-four extension.  With $W_0=1$ and
$W_3=W_1W_2/G$, the four $W_\chi$ form a basis.  At the
given integer point use positive square roots.  Then all
$W_\chi(x,y)$ are integers: $F_3(x,y)$ is an integer which
is a rational square, hence an integer square, and
$W_3(x,y)=\sqrt{F_1(x,y)}\sqrt{F_2(x,y)}/G(x,y)$.

\emph{The linear system at infinity.}
Put $r_0=s_0=0$, $k=3q$, and consider
\[
 \mathcal A=\sum_{\chi=0}^3A_\chi(X,Y)W_\chi,
 \qquad \deg_XA_\chi\le k-r_\chi,
 \quad \deg_YA_\chi\le k-s_\chi.
\]
There are
\[
 n=\sum_{\chi=0}^3(k-r_\chi+1)(k-s_\chi+1)
   \ge21q^2+18q+4
\]
unknown coefficients, and $n\le4(3q+1)^2$.
Writing $\xi=1/X$, $\eta=1/Y$, the positive formal branches
are
\[
 W_\chi=X^{r_\chi}Y^{s_\chi}
 \prod_{u\in I_\chi}(1+u\xi)^{1/2}
 \prod_{v\in J_\chi}(1+v\eta)^{1/2}.
\]
Thus $f(\xi,\eta)=\xi^k\eta^k\mathcal A$ is a formal
power series.  Require its coefficients of total degree at
most $2k$ to vanish.  This gives at most
\[
 t=\frac{(2k+1)(2k+2)}2=18q^2+9q+1,
 \qquad n-t\ge3q^2+9q+3,\qquad \frac t{n-t}<6
\]
homogeneous linear equations.

All coefficients required from the binomial series have
denominator dividing $4^{2k}$.  Cauchy's estimate on circles
of radius $(2B)^{-1}$ bounds their numerators, after this
common denominator is cleared, by $(B+1)^{c_0q}$ for an
absolute $c_0$.  An elementary pigeonhole form of Siegel's
lemma yields a nonzero integer solution of height
$(B+1)^{c_1q}$.  Indeed, if the integer matrix has height $A$,
the images of $\{0,\ldots,M\}^n$ take at most
$(2nAM+1)^t$ values.  Choosing
$M=\lceil(3nA)^7\rceil$ makes two images equal, because
$t/(n-t)<6$.  Their difference is a nonzero kernel vector;
$\log n=O(\log B)$ is absorbed in the stated height.

\emph{Primitive coefficients.}
Let $D$ be the primitive polynomial gcd of the four
$A_\chi$, omitting zero coefficients, and write
$A_\chi=DA'_\chi$.  The factor bound
of Lemma~\ref{p13:rig:factor-height} keeps
$H(A'_\chi)\le(B+1)^{O(q)}$.  Write
$a=\deg_XD$, $b=\deg_YD$, $e=\deg D$.  Then
\[
 f=\widetilde D f',\qquad
 \widetilde D=\xi^a\eta^bD(\xi^{-1},\eta^{-1}),\qquad
 f'=\xi^{k-a}\eta^{k-b}\sum_\chi A'_\chi W_\chi.
\]
These are power series, and
$\operatorname{ord}\widetilde D=a+b-e$ for total order.
Since order is additive in the power-series domain,
\begin{equation}
 \operatorname{ord}f'\ge2k+1-a-b+e
                          =(k-a)+(k-b)+e+1.
 \label{p13:rig:primitive-order}
\end{equation}
Thus removing $D$ preserves vanishing beyond the remaining
pole weight.  This is a formal division, never a division
by the possibly zero value $D(x,y)$.  Replace the auxiliary
by its primitive version from now on.

\emph{Integer separation and the norm.}
On the bidisc $|\xi|,|\eta|\le(4B)^{-1}$, the analytic
numerator is bounded by $(B+1)^{O(q)}$.  Cauchy's estimate
and the order of vanishing give, in the comparable range
$x,y\in[N,2N)$,
\begin{equation}
                  |\mathcal A(x,y)|\le
                                   \frac{(B+1)^{c_2q}}N.
 \label{p13:rig:integer-gap}
\end{equation}
To see the power of $N$ explicitly before removing $D$, the
tail of the numerator is bounded by its coefficient bound
times
\[
 \sum_{\ell\ge2k+1}(\ell+1)(4B/N)^\ell
       \ll(k+1)(4B/N)^{2k+1}.
\]
Multiplication by $x^ky^k\le(2N)^{2k}$ gives
\eqref{p13:rig:integer-gap}.  After removing $D$, both pole
weights decrease, and \eqref{p13:rig:primitive-order} gives
the same estimate with the new height bound.

The specialization $\mathcal A(x,y)$ is an integer, so it
vanishes for the stated threshold.  The auxiliary is nonzero
in $K$ by independence of the four characters.  Its norm is
therefore nonzero.  An explicit polynomial expression is
\[
 \begin{split}
 U&=A_0^2+A_1^2F_1-A_2^2F_2-A_3^2F_3,\\
 V&=A_0A_1-A_2A_3(F_2/G),\\
 R&=U^2-4F_1V^2.
 \end{split}
\]
Here $F_2/G$ is a polynomial.  The partial degrees of $U$
are at most $2k$, and those of $V$ are at most
$2k-r_1$, $2k-s_1$.  For the second term in $V$, use
$r_3=r_1+r_2-\deg_XG$ and its $Y$ analogue.
It follows that the partial degrees of $R$ are at most
$4k=12q$.  The coefficient $\ell^1$ norms of all factors
in this formula are $(B+1)^{O(q)}$, which proves the height
bound and completes \eqref{p13:rig:two-polynomial}.
\end{proof}


\subsection{Proof of Theorem~\ref{p13:rig:common-curve}}
\label{p14:comp:rigidity:common-curve}

\begin{proof}
For $i<j$, the half-cardinalities of $w_i,w_j,w_i+w_j$
are at most $q=\min(2q_0,B)$, since an even subset of
$[-B,B]\cap\Z$ has size at most $2B$.
Proposition~\ref{p13:rig:two-witnesses} gives a nonzero
polynomial $R_{ij}$ of degree at most $48q_0$ and height
$(B+1)^{O(q_0)}$, vanishing at $(x,y)$ and arising from
an auxiliary with primitive coefficients.
Choose an irreducible primitive factor $C_0$ of $R_{12}$
which vanishes at this point.  The factor-height lemma
preserves its height bound.  If $C_0$ depended on only one
variable, its integer root of size at least $N$ would
contradict Cauchy's root bound.  Thus it depends on both.

If $C_0$ failed to divide some $R_{ij}$, the resultant
$\operatorname{Res}_Y(C_0,R_{ij})$ would be a nonzero integer
polynomial vanishing at $x$.  Its height is
$(B+1)^{O(q_0^2)}$: its Sylvester matrix has size $O(q_0)$,
entries of coefficient norm $(B+1)^{O(q_0)}$, and
$O(q_0)!$ determinant terms.  This again contradicts the
root bound under the stated threshold.  Hence $C_0$ divides
every $R_{ij}$.

No factor $X+u$ or $Y+v$ vanishes identically on this curve,
since $C_0(x,y)=0$ with both coordinates larger than $B$.
The displayed square classes are therefore defined in
$K_0$.  If the classes of $F_{w_i},F_{w_j}$ were independent
there, they would still generate a degree-four field
extension.  Specializing the polynomial norm formula gives
zero, since $C_0\mid R_{ij}$; in a field this forces the
auxiliary to vanish.  Independence of its four characters
would force all four coefficient polynomials to be divisible
by $C_0$, contrary to their primitivity.  Thus the images
of every pair of basis vectors are dependent.  Over $\F_2$
this means that all nonzero images are equal, so the image
has dimension at most one.
\end{proof}

\subsection{Proof of Proposition~\ref{p13:rig:aligned}}
\label{p14:comp:rigidity:aligned}

\begin{proof}
Fix $\tau<1$.  Theorem~\ref{p13:rig:short-basis} gives a
space $V=\mathcal K_\tau$ of dimension
\[
 t\ge d_{xy}-\frac23B\log(1/\tau)
                     -O_{\tau,C}(B/\sqrt{\log B}),
\]
with a basis of half-cardinalities at most
$q_0=B^{\tau+o(1)}$.  The case $t<2$ is covered by
the displayed error term.  For $t\ge2$, apply
Proposition~\ref{p13:rig:two-witnesses} to each pair of
basis vectors.  Its threshold holds because
$q_0\log(B+1)=o(B)=o(\log N)$.

Put $\ell(X,Y)=bX-aY-h$.  Every resulting norm polynomial
$R_{ij}$ is divisible by $\ell$.  Otherwise, with
$d_Y=\deg_YR_{ij}$, the nonzero integer polynomial
\[
           a^{d_Y}R_{ij}\left(X,\frac{bX-h}{a}\right)
\]
would vanish at $x$.  Its height is $(B+1)^{O(q_0)}$
since $\max(a,b,|h|)\le2B^4$, contradicting the
integer root bound for $x\ge N$.
The primitive-norm argument in
Theorem~\ref{p13:rig:common-curve} now applies to this
line and shows that the image of $V$ in
$\mathbb Q(\ell=0)^\times/\mathbb Q(\ell=0)^{\times2}$
has dimension at most one.

Let $V_0$ be its kernel, of dimension at least $t-1$.
On the line, $X+u$ and $Y+v$ have the same root precisely when
\begin{equation}
                            h+bu=av.
 \label{p13:rig:line-matching}
\end{equation}
Every word of $V_0$ must pair all its first-block roots
with second-block roots.  The integer solutions of
\eqref{p13:rig:line-matching} advance by $(a,b)$ in two
intervals of $B$ consecutive offsets, so their number $r$
satisfies
\[
                        r\le1+\frac{B-1}{\max(a,b)}.
\]
Injectivity of the physical projection gives
$t-1\le\dim V_0\le r$, proving
\eqref{p13:rig:aligned-quantitative} after absorbing the
bounded terms.

For compatible distinct centres with $\max(a,b)\ge2$,
this gives $d_{xy}\le B/2+(2/3)B\log(1/\tau)+o_{\tau,C}(B)$.
Choose $\tau$ sufficiently close to one after fixing a
desired error margin, then let $N$ tend to infinity.
This order of choices proves the uniform $B/2+o_C(B)$ bound.

For $a=b=1$, compatibility implies $|x-y|\ge L$:
if $x<y<x+L$, the first run makes $f(y-1)=f(y)$,
whereas the second requires their signs to be opposite.
At distance exactly $L=B-1$, the true windows share one
position; at larger distance they are disjoint.  Thus
\eqref{p13:rig:line-matching} has at most one solution.
A nonzero word in $V_0$ would need at least two pairings,
because its selections are nonempty and even.  Consequently
$V_0=0$, $t\le1$, and the dimension estimate followed by
$\tau\uparrow1$ gives $d_{xy}=o_C(B)$.
\end{proof}
